\documentclass[10pt,a4paper]{amsart}
\usepackage[margin=19mm]{geometry}
\usepackage{amsmath,amssymb,mathtools}
\usepackage[scr=rsfs,cal=euler]{mathalfa}
\usepackage{xfrac}
\usepackage[unicode,hidelinks,bookmarks=false]{hyperref}
\newcommand{\ie}{i.e.\ }
\newcommand{\cf}{cf.\ }
\newcommand{\resp}{respectively\ }
\newcommand{\Subsection}{\subsection}
\providecommand{\nobreakdash}{\nobreak\hbox{-}}
\setlength{\emergencystretch}{2em}
\multlinegap0pt
\usepackage{amssymb}
\usepackage{xcolor}
\newtheorem{proposition}{Proposition}
\newtheorem{claim}{Claim}
\newcommand{\bfs}{\boldsymbol}
\def\ifm#1#2{\relax \ifmmode#1\else#2\fi}
\newcommand{\klk}    {\ifm {,\ldots,} {$,\ldots,$}}
\title{A Kronecker algorithm for locally closed sets over a perfect field}
\author{Nardo Gim\'enez, Joos Heintz, Guillermo Matera, Luis Miguel Pardo, Mariana P\'erez, Melina Privitelli}
\date{Source: arXiv:2512.14888v2 (CC BY 4.0)}
\begin{document}
\maketitle

\noindent\emph{Source and licence.} A Kronecker algorithm for locally closed sets over a perfect field. Version arXiv:2512.14888v2 (CC BY 4.0), distributed by arXiv under the Creative Commons Attribution 4.0 International licence, http://creativecommons.org/licenses/by/4.0/, as recorded in the arXiv licence metadata for this exact version and shown on the abstract page https://arxiv.org/abs/2512.14888v2. Full licence notice: \texttt{licenses/arxiv-2512.14888-CC-BY-4.0.txt}.\par\smallskip

\noindent\textit{Editorial scope.} Counted unit: the article's proposition lifting (source0.tex lines 4162--4197). The article's complete original proof of it is delivered with it (source0.tex lines 4199--4410, one \texttt{proof} environment of 212 lines). No mathematics is written, repaired or re-derived: the statement and the proof are the article's own lines, in the article's own notation, and no retained line is substituted or re-flowed.  No further source-local context is needed: every cross-reference of every retained line resolves inside the two delivered spans.  Nothing was omitted from any retained span, so the package carries no omission record.  No retained line contains a citation, so the wrapper carries no bibliography and no bibliography entry.  No retained line contains a figure environment, an image-inclusion command or drawing code, so no image is delivered and the package declares no asset dependency.  Editorial formatting only, in addition: an AMS \texttt{amsart} preamble that restates the article's own declarations and macros used by the retained lines, namely the environments \texttt{proposition}, \texttt{claim} on separate counters, together with the standard packages those lines need.  The retained environments print in this document as Proposition 1, Claim 1 in order of appearance, whereas the article prints the same items with the numbering of its own full document.  The complete native source of the article as distributed in the arXiv e-print for this exact version is bundled unchanged as \texttt{sources/191127/source0.tex}.
\begin{proposition}\label{prop: main loop Newton lifting}
Fix $j\ge 0$, and assume that we are given polynomials $M^j,V_{n-s+1}^j, \ldots V_n^j\in
k[\![X_{n-s} -p_{n-s}]\!][T]$, with $\deg_TM^j=\delta_s$ and $\deg_T
V_{n-s+i}^j<\delta_s$ for $1\le i\le s$, satisfying the following
congruences with precision $2^j$:
%
\begin{align}
F_i (\bfs p^{s+1}, X_{n-s}, \bfs V^j) & \equiv 0  \mod M^j\quad (1 \le i
\le s),
\label{eq: CongruenceProofHL1}\\
V_{n-s+1}^j& = T, \label{eq: CongruenceProofHL2}
\end{align}
%
where $\bfs V^j:=(V_{n-s+1}^j, \ldots, V_n^j)$. Assume further that
$\det J_s(X_{n-s}, \bfs V^j)$ is invertible modulo $M^j$. Then the
$(j+1)$th iteration of the loop of step {\bf 4} can be correctly
executed and outputs polynomials $M^{j+1},V_{n-s+1}^{j+1}, \ldots
V_n^{j+1} \in k[\![X_{n-s} -p_{n-s}]\!][T]$, with
$\deg_TM^{j+1}=\delta_s$ and $\deg_T V_{n-s+i}^{j+1}<\delta_s$ for
$1\le i\le s$, satisfying the following congruences with precision
$2^{j+1}$:
%
\begin{align*}
F_i (\bfs p^{s+1}, X_{n-s},\bfs V^{j+1})  & \equiv 0  \mod M^{j+1} \quad
(1 \le i \le s),
\\
V_{n-s+1}^{j+1}& = T,
\end{align*}
%
where $\bfs V^{j+1}\!:=(V_{n-s+1}^{j+1}, \ldots, V_n^{j+1})$.
Moreover, $\det J_s(X_{n-s},\bfs V^{j+1} )$ is invertible modulo
$M^{j+1}$. In addition, if $G(\bfs p^{s+1}, X_{n-s}, \bfs V^j)\in
k[\![X_{n-s} -p_{n-s}]\!][T]$ is invertible modulo $M^j$, then
$G(\bfs p^{s+1}, X_{n-s}, \bfs V^{j+1})$ is invertible modulo
$M^{j+1}$.
\end{proposition}

\begin{proof}
Regarding the execution of the $(j+1)$th iteration, the only point
to check is whether the Newton-Hensel operator $N_{\bfs F_s}
(X_{n-s}, \bfs V^j)$ is well-defined. By hypothesis, $\det J_s(
X_{n-s},\bfs V^j)\mod M^j$ is invertible, and thus the
matrix $J_s(X_{n-s},\bfs V^j)$ admits an inverse modulo
$M^j$. Therefore, $N_{\bfs F_s} (X_{n-s}, \bfs V^j)$
is well--defined modulo $M^j$.

Now we analyze the output produced by the $(j+1)$th iteration. Set $\bfs
v^{j+1}:=(v_{n-s+1}^{j+1}\klk v_n^{j+1})$. We first claim that
%
\begin{equation}\label{eq: F_v(j+1)_=_0_mod_M_j}
\bfs F_s(X_{n-s}, \bfs v^{j+1}) \mod M^j = 0
\end{equation}
%
with precision $2^{j+1}$. Indeed, from (\ref{eq: CongruenceProofHL1})
we have $\bfs F_s (X_{n-s},\bfs V^j) = 0\mod M^j$ with precision $2^j$.
By the definition of $\bfs v^{j+1}$, it follows that
%
$$\bfs v^{j+1}-\bfs V^j = -
(J_s^{-1} \bfs F_s) (X_{n-s},\bfs V^j)
=0\mod M^j,$$
%
with precision
$2^j$. Since
$\deg_T(v_{n-s+i}^{j+1}-V_{n-s+i}^j)<\delta_s$ for $1\le i\le s$, we
conclude that
%
\begin{equation}\label{eq: Congruence w(j+1) W(j)}
\bfs v^{j+1}-\bfs V^j = 0
\end{equation}
%
with precision $2^j$, and thus, $(\bfs v^{j+1}-\bfs V^j)^2 = 0$ with
precision $2^{j+1}$. Applying the second-order Taylor expansion of
$\bfs F_s$ in powers of $X_{n-s+1}- V_{n-s+1}^j\klk X_n - V_n^j$, we
obtain
%
\begin{align*}
\bfs F_s(X_{n-s}, \bfs v^{j+1}) & \equiv  \bfs F_s(X_{n-s},\bfs V^j)+J_s(X_{n-s}, \bfs V^j)
(\bfs v^{j+1}-\bfs V^j) \mod M^j \\
\nonumber  & \equiv \bfs F_s(X_{n-s, }\bfs V^j) - J_s(X_{n-s},\bfs
V^j)(J_s^{-1} \bfs F_s) (X_{n-s},\bfs V^j)\mod M^j=\bfs 0,
\end{align*}
%
with precision $2^{j+1}$. This proves the claim. Note that, as
$\bfs v^{j+1}$ agrees with $\bfs V^j$ with precision $2^j$, step
{\bf 4.1} in the main loop only involves modifications of the
coefficients of $\bfs V^j$ of order $2^j$ or higher.

It follows that $\bfs F_s$ vanishes at
$\bfs v^{j+1}$ with precision $2^{j+1}$, but modulo $M^j$,
which represents the output with
precision only $2^j$. The next steps in the main loop
correct this lack of precision. To analyze these steps, recall that
we define
%
$$\Delta^{j+1}:=v_{n-s+1}^{j+1}- T.$$
%
We may consider $\Delta^{j+1}$ as the ``error'' or ``failure'' of
$\bfs v^{j+1}$ to correctly parametrize the coordinates of
$X_{n-s+1}\klk X_n$ with precision $2^{j+1}$. Observe that
%
$$\Delta^{j+1}=0,\quad \Delta^{j+1}_{M}=0,\quad \mathrm{and}\quad
\Delta_i^{j+1}=0\quad (n-s+1 \le i \le n),$$
%
all with precision ${2^j}$. In fact, $\Delta^{j+1}=0$ with precision
${2^j}$ because $v_{n-s+i}^{j+1}=V_{n-s+i}^j$ with precision $2^j$
and (\ref{eq: CongruenceProofHL2}) holds, while the remaining
assertions follow immediately from this one. The fact that
$\Delta^{j+1}=0$ with precision $2^j$ indicates that $\bfs v^{j+1}$
correctly parametrizes $X_{n-s+1}\klk X_n$ with precision $2^j$.

\begin{claim} Fix $j\geq 0$. The following congruences hold with precision $2^{j+1}$:
 \begin{align}
% \Delta^{j+1}&=0  \ \ \textrm{with precision} \ \ 2^j,\label{eq: lemma: main loop Newton lifting_1}\\
% \Delta^{j+1}(T\pm\Delta^{j+1})&=\Delta^{j+1}, \label{eq: lemma: main loop Newton lifting_2}\\
 M^j(T-\Delta^{j+1})&\equiv 0 \mod M^{j+1},\label{eq: lemma: main loop Newton lifting_3}\\
 v^{j+1}_i(T-\Delta^{j+1})&\equiv V^{j+1}_i \mod M^{j+1} \textrm{ for }n-s+1\le i \le n.
 \label{eq: lemma: main loop Newton lifting_4}
 \end{align}
\end{claim}
%
\begin{proof}[Proof of Claim]
%The Taylor expansion of $\Delta^{j+1}(T\pm\Delta^{j+1})$ in powers
%of $\Delta^{j+1}$ is
%$$
%\Delta^{j+1}(T\pm\Delta^{j+1})=\Delta^{j+1}\pm \frac{\partial
%\Delta^{j+1}}{\partial T} \Delta^{j+1} +
%\mathcal{O}\big((\Delta^{j+1})^2\big).
%$$
%Since $\Delta^{j+1}=0$ with precision $2^{j}$, we have
%$\frac{\partial \Delta^{j+1}}{\partial T}\Delta^{j+1}=0$ and
%$(\Delta^{j+1})^2=0$ with precision $2^{j+1}$. This proves
%\eqref{eq: lemma: main loop Newton lifting_2}.
%
We start by expanding $M^j(T-\Delta^{j+1})$ using Taylor expansion
around $T$, in
powers of $\Delta^{j+1}$:
%
$$
M^j(T-\Delta^{j+1})=M^j(T)-\frac{\partial M^j}{\partial T}\Delta^{j+1}
+\mathcal{O}\big((\Delta^{j+1})^2\big).
$$
%
Since $\Delta^{j+1}=0$ with precision $2^{j}$, it follows that
$(\Delta^{j+1})^2=0$ with precision $2^{j+1}$. Hence,
%
\begin{equation}\label{eq: Taylor expansion M}
M^j(T-\Delta^{j+1})=M^j(T)-\frac{\partial M^j}{\partial T}\Delta^{j+1}
+\mathcal{O}\big((\Delta^{j+1})^2\big),
\end{equation}
%
with precision $2^{j+1}$. By the definition of $M^{j+1}$,
we have $M^j=M^{j+1}+\Delta^{j+1}_{M}$, where $\Delta^{j+1}_{M}$ is the
remainder of the division of $\Delta^{j+1}\frac{\partial
M^j}{\partial T}$ by $M^j$, taken with precision $2^{j+1}$. Let
$h\in k[\![X_{n-s}-p_{n-s}]\!][T]$ denote the quotient of this
division. By \eqref{eq: Taylor expansion M}, and the definition
of $M^{j+1}$ and $\Delta^{j+1}_{M}$, we have
%
\begin{align*}
M^j(T-\Delta^{j+1})=
M^j-(h(T)M^j+\Delta^{j+1}_{M})
&=
(M^{j+1}+\Delta^{j+1}_{M})\big(1-h(T)\big)-\Delta^{j+1}_{M}\\
&=M^{j+1}\big(1-h(T)\big)- h(T)\Delta^{j+1}_{M}(T),
\end{align*}
%
with precision $2^{j+1}$. Since $\Delta^{j+1}\frac{\partial
M^j}{\partial T}=0$ with precision $2^{j}$, and $M^j$ is monic,
the quotient $h(T)$ vanishes with precision $2^{j}$, and thus
$\Delta^{j+1}_{M}=0$ with precision $2^{j}$. We conclude that
$h(T)\Delta^{j+1}_{M}(T)=0$ with precision $2^{j+1}$, which implies
\eqref{eq: lemma: main loop Newton lifting_3}.

Now we turn to \eqref{eq: lemma: main loop Newton
lifting_4}. For $n-s+1\le i\le n$, consider the Taylor expansion of
$v^{j+1}_i(T-\Delta^{j+1})$ around $T$, in
powers of $\Delta^{j+1}$. We have
$$
v^{j+1}_i(T-\Delta^{j+1})=v^{j+1}_i(T)-\Delta^{j+1}\frac{\partial v^{j+1}_i}{\partial T}(T),
$$
with precision $2^{j+1}$. By construction,
$v^{j+1}_i=V^{j+1}_i+\Delta^{j+1}_i$ with precision $2^{j+1}$, where
$\Delta^{j+1}_{i}$ is the remainder of the division of
$\Delta^{j+1}\frac{\partial v^{j+1}_i}{\partial T}$ by $M^j$, taken
with precision $2^{j+1}$. Let $h_i\in k[\![X_{n-s}-p_{n-s}]\!][T]$
denote the quotient of this division. Then
%
\begin{equation}\label{eq: proof: main loop Newton lifting_2}v^{j+1}_i(T-\Delta^{j+1})=V_i^{j+1}+\Delta^{j+1}_i-(h_i(T)M^j(T)+\Delta_i^{j+1})
=V_i^{j+1}-h_i(T)M^j(T),
\end{equation}
%
with precision $2^{j+1}$. We observe that $h_i(T)M^j(T)\equiv 0\mod
M^{j+1}$ with precision $2^{j+1}$. Indeed, similarly as before we
see that $h_i=0$, and then $\frac{\partial h_i}{\partial T}=0$, with
precision $2^{j}$. It follows that $
h_i(T-\Delta^{j+1})=h_i(T)-\frac{\partial h_i}{\partial
T}\Delta^{j+1} + \mathcal{O}\big((\Delta^{j+1})^2\big)=h_i(T) $ with
precision $2^{j+1}$. Then, by \eqref{eq: lemma: main loop Newton
lifting_3},
%
\begin{align*}
h_i(T)M^j(T)=h_i(T)\Big(M^j(T)-\frac{\partial M^j}{\partial T}\Delta^{j+1}\Big)&=
h_i(T-\Delta^{j+1})M^j(T-\Delta^{j+1})\\&\equiv 0\mod M^{j+1},
\end{align*}
%
with precision $2^{j+1}$. This, combined with \eqref{eq: proof: main
loop Newton lifting_2}, yields \eqref{eq: lemma: main loop Newton
lifting_4}, which completes the proof of the claim.
\end{proof}

Substituting $T-\Delta^{j+1}$ for $T$ in
\eqref{eq: F_v(j+1)_=_0_mod_M_j}, we obtain
%
\begin{equation}\label{eq: F_v(j+1)_=_0_mod_M_j_1}
\bfs F_s\big(X_{n-s}, \bfs v^{j+1}(T-\Delta^{j+1})\big)\equiv 0 \mod M^{j}(T-\Delta^{j+1}),
\end{equation}
%
with precision $2^{j+1}$. Taking into account \eqref{eq: lemma: main
loop Newton lifting_3} and \eqref{eq: lemma: main loop Newton
lifting_4}, we conclude that $$\bfs F_s\big(X_{n-s}, \bfs V^{j+1}
\big) \equiv 0 \mod M^{j+1}$$ with precision $2^{j+1}$, as desired.
Furthermore, since by definition $v^{j+1}_{n-s+1}=T+\Delta^{j+1}$,
substituting $T-\Delta^{j+1}$ for $T$ in this expression we obtain
$V^{j+1}_{n-s+1}\equiv T \mod M^{j+1}$, with precision $2^{j+1}$.

It remains to prove the invertibility of $\det J_s( X_{n-s}, \bfs
V^{j+1})$ and $G(\bfs p^{s+1}, X_{n-s}, \bfs V^{j+1})$, both modulo
$M^{j+1}$. Observe that the invertibility of the former holds if and
only if the resultant $\mathrm{Res}_T\big(\det J_s( X_{n-s}, \bfs
V^{j+1}),M^{j+1}\big)$ is an invertible element of
$\overline{k}[\![X_{n-s} -p_{n-s}]\!]$. By \eqref{eq: Congruence
w(j+1) W(j)}, we have $\bfs v^{j+1}=\bfs V^j$ with precision $2^j$.
Furthermore, as $V_i^{j+1}:=v_i^{j+1}- \Delta_i^{j+1}$, and
$\Delta_i^{j+1}=0$ with precision $2^j$, for $1\le i\le s$, it
follows that $\bfs V^{j+1}=\bfs V^j$ with precision $2^j$, which in
turn implies $\det J_s( X_{n-s}, \bfs V^{j+1})=\det J_s(X_{n-s},
\bfs V^j)$ with precision $2^j$. As $M^{j+1}=M^j$ with precision
$2^j$, we conclude that
%
$$\mathrm{Res}_T\big(\det J_s(X_{n-s}, \bfs
V^{j+1}),M^{j+1}\big)=\mathrm{Res}_T\big(\det J_s(
X_{n-s}, \bfs V^j),M^j\big)$$
%
with precision $2^j$. Since the right-hand side is invertible by
hypothesis, the assertion follows. The same argument proves that the
invertibility of $G(\bfs p^{s+1}, X_{n-s}, \bfs V^j)$ modulo $M^j$
implies the invertibility of $G(\bfs p^{s+1}, X_{n-s}, \bfs V^{j+1})$
modulo $M^{j+1}$.
\end{proof}

\end{document}
